% =============================================================================
% Chapter 7: Conclusion and Future Work
% =============================================================================

\chapter{Conclusion and Future Work}
\label{ch:conclusion}


% =============================================================================
\section{Conclusion}
\label{sec:final-conclusion}
% =============================================================================

This thesis addressed the problem of file-level code retrieval for AI coding agents: given a natural-language task description and a large Python repository, which source files should the agent examine? We developed CodeGraph, a hybrid retrieval system that combines lexical seed extraction with Personalized PageRank over a deterministic code dependency graph, and evaluated it on 300 real-world issues from the SWE-bench Lite benchmark.

\subsection{Contributions}

We make three primary contributions:

\begin{enumerate}
\item \textbf{The CodeGraph Retrieval Pipeline.} We designed and implemented a hybrid architecture that transforms Python repositories into deterministic code dependency graphs. By combining multi-signal seed selection with Personalized PageRank, the system provides a ranked and inspectable view of a repository designed for AI agent use.

\item \textbf{Evaluation Methodology and Failure Taxonomy.} We evaluated the system on SWE-bench Lite, achieving 74.0\% Recall@10 and demonstrating a 24.3 percentage point gain over lexical baselines. By analyzing both successful and failed retrievals, we identified reachability as the primary retrieval bottleneck, leading to a seed-first engineering strategy. We classified the remaining failures into a taxonomy of connectivity gaps and semantic gaps, defining the structural ceiling of graph-based retrieval.

\item \textbf{Deterministic Agent Integration.} We implemented a stable, MCP-based interface that exposes structural retrieval to AI agents. This interface provides not only ranked context but also seed provenance and structural paths, making the retrieval process transparent and reproducible for both human developers and autonomous agents.

\end{enumerate}

\subsection{Central Finding}

The central finding of this thesis is that lexical matching provides reachability, while graph propagation provides ranking. Textual similarity can locate the general region of a codebase where a task is relevant, but it cannot navigate the dependency chains that connect related files. Graph propagation follows these chains, using the repository's structure to rank files by their functional relationship to the initial match. This confirms that codebases are connected systems, not collections of independent documents, and that retrieval methods must account for this structure.

% =============================================================================
\section{Future Work}
\label{sec:future-work}
% =============================================================================

The failure taxonomy identifies the most productive directions for future research, centered on bridging the gap between natural language and codebase structure. Vector embeddings could connect natural-language descriptions to code entities through learned semantic similarity, providing seeds that BM25 cannot reach. Alternatively, a language model could preprocess issue descriptions to suggest likely module or class names, augmenting the lexical seeds with concrete entity suggestions.

To close connectivity gaps, future work could explore selective sibling linkages within cohesive directories or gate new edges on co-change evidence from version control. Additionally, type-informed call resolution using lightweight type inference from annotations or property-chain analysis could add edges to the graph without requiring full runtime information.

The pipeline's language-agnostic design allows for extension to TypeScript or Java through new parser and resolution logic. The next step is to measure the end-to-end impact: how often does an agent produce a correct patch with structural context versus without it? Plugging CodeGraph into a complete SWE-bench resolution pipeline would test this link directly.

This thesis demonstrates that codebases are connected systems, and that retrieval methods must account for this structure. By combining lexical matching with graph propagation, CodeGraph provides a deterministic foundation for repository-scale code navigation. As AI coding agents continue to take on more complex tasks, the integration of structural context will be essential to their effectiveness.
